\begin{derivation}[从\eqnrefs{eq:feshbach-block-matrix-dp68}到\eqnrefs{eq:feshbach-schur-proof-r10}]

正文\eqnrefs{eq:feshbach-block-matrix-dp68} 把完整预解式的逆算符分成目标块 $P$ 和补块 $Q$。只要 $z-QHQ$ 可逆，就能先从下方分块方程消去 $Q$ 分量，再代回 $P$ 方程；这一普通分块求逆步骤直接给出正文\eqnrefs{eq:feshbach-schur-proof-r10}，不需要弱耦合假设。

记
\begin{equation*}
 z-H=
 \begin{pmatrix}A&B\\C&D\end{pmatrix},
 \quad
 A=z-PHP,
 \quad B=-PHQ,
 \quad C=-QHP,
 \quad D=z-QHQ.
\end{equation*}
设其逆矩阵为
\begin{equation*}
 \begin{pmatrix}X&Y\\Z&W\end{pmatrix}.
\end{equation*}
由下左块方程
\begin{equation*}
 CX+DZ=0
\end{equation*}
得到
\begin{equation*}
 Z=-D^{-1}CX.
\end{equation*}
代入上左块 $AX+BZ=I_P$，
\begin{align*}
 (A-BD^{-1}C)X&=I_P,\\
 X&=(A-BD^{-1}C)^{-1}.
\end{align*}
而 $X=P(z-H)^{-1}P$。恢复四个分块以后即得正文\eqnrefs{eq:feshbach-schur-proof-r10}。
\end{derivation}

\begin{derivation}[从\eqnrefs{eq:sw-sylvester-r15}到\eqnrefs{eq:sw_arbitrary_target_subspace_matrix_elements}]
Baker--Campbell--Hausdorff 展开给出

\begin{align*}
 e^SHe^{-S}
 &=H+[S,H]+\frac12[S,[S,H]]+\order(S^3H)\\
 &=H_{\rm d}+H_{\rm od}+[S,H_{\rm d}]+[S,H_{\rm od}]
 +\frac12[S,[S,H_{\rm d}]]+\order(S^3H).
\end{align*}
由正文\eqnrefs{eq:sw-sylvester-r15}，$[S,H_{\rm d}]=-H_{\rm od}$，因此
\begin{align*}
 e^SHe^{-S}
 &=H_{\rm d}+[S,H_{\rm od}]
 -\frac12[S,H_{\rm od}]+\order(S^3H)\\
 &=H_{\rm d}+\frac12[S,H_{\rm od}]+\order(S^3H),
\end{align*}
即正文\eqnrefs{eq:sw-bch-second-r15}。

令
\begin{equation*}
 H_{\rm d}|p\rangle=E_p|p\rangle,
 \qquad
 H_{\rm d}|q\rangle=E_q|q\rangle,
\end{equation*}
其中 $p,p'\in P$、$q\in Q$。在 $\langle q|\cdots|p\rangle$ 上取正文\eqnrefs{eq:sw-sylvester-r15} 的矩阵元，
\begin{align*}
 0&=H_{qp}+(E_p-E_q)S_{qp},\\
 S_{qp}&=\frac{H_{qp}}{E_q-E_p},
 \qquad
 S_{pq}=\frac{H_{pq}}{E_p-E_q}.
\end{align*}
于是
\begin{align*}
 \langle p|[S,H_{\rm od}]|p'\rangle
 &=\sum_{q\in Q}
 \left(S_{pq}H_{qp'}-H_{pq}S_{qp'}\right)\\
 &=\sum_{q\in Q}H_{pq}H_{qp'}
 \left[
 \frac1{E_p-E_q}
 +\frac1{E_{p'}-E_q}
 \right].
\end{align*}
乘上正文\eqnrefs{eq:sw-bch-second-r15} 中的 $1/2$，并加回 $P H_{\rm d}P$，得到正文\eqnrefs{eq:sw_arbitrary_target_subspace_matrix_elements}。
\end{derivation}

\begin{derivation}[从\eqnrefs{eq:resolvent-spectral-resolution-dp68}到\eqnrefs{eq:resolvent-norm-spectral-dp21}]

正文\eqnrefs{eq:resolvent-spectral-resolution-dp68} 把自伴算符的预解式写成谱测度积分。对任意归一化态取范数平方，可以把算符范数问题化成正概率测度下的标量函数 $|z-\lambda|^{-2}$ 的上确界，从而得到正文\eqnrefs{eq:resolvent-norm-spectral-dp21}。

对归一化 $|\psi\rangle$ 定义
\begin{equation*}
 \dd\mu_\psi(\lambda)
 =\langle\psi|\dd E_A(\lambda)|\psi\rangle,
 \qquad
 \int\dd\mu_\psi=1.
\end{equation*}
于是
\begin{align*}
 \|(z-A)^{-1}|\psi\rangle\|^2
 &=\int_{\sigma(A)}
 \frac{\dd\mu_\psi(\lambda)}{|z-\lambda|^2}\\
 &\le
 \sup_{\lambda\in\sigma(A)}\frac1{|z-\lambda|^2}.
\end{align*}
对所有归一化态取上确界，得到
\begin{equation*}
 \|(z-A)^{-1}\|
 \le\frac1{\operatorname{dist}[z,\sigma(A)]}.
\end{equation*}
反向不等式可取一列谱权重越来越集中在离 $z$ 最近谱邻域的归一化态；离散谱情形就是取最近本征态。因此等号成立。令 $A=QHQ$ 即得正文\eqnrefs{eq:resolvent-norm-spectral-dp21}。
\end{derivation}

\begin{derivation}[从\eqnrefs{eq:feshbach_schur_exact_positive_projection}到\eqnrefs{eq:feshbach-second-order-static-dp68}]

精确 Feshbach--Schur 有效算符为
\[
H_{\rm eff}(z)
=
PHP+\Sigma_P(z),
\qquad
\Sigma_P(z)
=
PHQ\,(z-QHQ)^{-1}QHP.
\]
这里 \(P\) 可以是任意有限维目标子空间，并不要求只保留一个本征态。设目标能窗以参考能量 \(E_*\) 为中心，并与 \(QHQ\) 的谱保持最小距离
\[
\Delta_Q
\equiv
\operatorname{dist}\!\left(E_*,\sigma(QHQ)\right)>0.
\]
若实际极点满足 \(|z-E_*|<\Delta_Q\)，则
\[
z-QHQ
=
(E_*-QHQ)
\left[
I+(z-E_*)(E_*-QHQ)^{-1}
\right].
\]
定义
\[
X\equiv
(z-E_*)(E_*-QHQ)^{-1}.
\]
由预解式范数
\[
\|(E_*-QHQ)^{-1}\|
\le\frac1{\Delta_Q}
\]
可知
\[
\|X\|
\le\frac{|z-E_*|}{\Delta_Q}<1,
\]
所以在目标能窗内可以作算符 Neumann 展开：
\begin{align*}
(z-QHQ)^{-1}
&=
\left(I+X\right)^{-1}(E_*-QHQ)^{-1}\\
&=
\left(I-X+X^2-\cdots\right)
(E_*-QHQ)^{-1}\\
&=
(E_*-QHQ)^{-1}
-(z-E_*)(E_*-QHQ)^{-2}
+\mathcal O\!\left(
\frac{|z-E_*|^2}{\Delta_Q^3}
\right).
\end{align*}
因此
\begin{align*}
\Sigma_P(z)
&=
PHQ(E_*-QHQ)^{-1}QHP\\
&\quad
-(z-E_*)
PHQ(E_*-QHQ)^{-2}QHP
+\cdots .
\end{align*}

若 \(V\equiv PHQ+QHP\) 是弱混合，并定义
\[
\epsilon_{\rm mix}
\equiv
\frac{\|PHQ\|}{\Delta_Q}\ll1,
\]
则第一项满足
\[
\|\Sigma_P(E_*)\|
\le
\frac{\|PHQ\|^2}{\Delta_Q}
=
\epsilon_{\rm mix}^2\Delta_Q.
\]
精确极点相对未耦合目标能量的移动于是也是二阶，
\[
|z-E_*|
=
\mathcal O(\epsilon_{\rm mix}^2\Delta_Q).
\]
代回第二项可得
\[
\left\|
(z-E_*)
PHQ(E_*-QHQ)^{-2}QHP
\right\|
=
\mathcal O(\epsilon_{\rm mix}^4\Delta_Q),
\]
所以在二阶精度下，自能中的能量依赖可以冻结在 \(E_*\)：
\[
H_{\rm eff}^{(2)}
=
PHP
+
PHQ(E_*-QHQ)^{-1}QHP
+
\mathcal O(\epsilon_{\rm mix}^4\Delta_Q).
\]

正文\eqnrefs{eq:feshbach-second-order-static-dp68} 是一维目标子空间的特例。取
\[
P=|p\rangle\langle p|,
\qquad
PHP=E_pP,
\qquad
E_*=E_p,
\]
并对 \(QHQ\) 插入本征态完备关系
\[
Q
=
\sum_{q\in Q}|q\rangle\langle q|,
\qquad
QHQ|q\rangle=E_q|q\rangle,
\]
得到
\begin{align*}
\langle p|\Sigma_P(E_p)|p\rangle
&=
\sum_{q\in Q}
\frac{
\langle p|H|q\rangle
\langle q|H|p\rangle
}{E_p-E_q}\\
&=
\sum_{q\in Q}
\frac{|H_{pq}|^2}{E_p-E_q}.
\end{align*}
于是
\[
z
=
E_p
+
\sum_{q\in Q}
\frac{|H_{pq}|^2}{E_p-E_q}
+
\mathcal O(\epsilon_{\rm mix}^4\Delta_Q),
\]
即正文\eqnrefs{eq:feshbach-second-order-static-dp68}。因此“二阶能量修正”不是从非简并微扰论类比得到，而是精确 Schur 补在谱隙与弱混合控制下的静态展开；若目标子空间简并或近简并，应保留整个 \(P\) 块而不能先选单个 \(|p\rangle\)。
\end{derivation}
